% Figure: classification of second-order linear PDEs by discriminant.
% Schematic placing the three canonical equations in the
% (A, C) plane of the operator A u_xx + 2B u_xy + C u_yy, with the
% sign of the discriminant B^2 - AC fixing the type.
\begin{figure}[htbp]
\centering
\begin{tikzpicture}
\begin{axis}[
  width=11cm, height=7cm,
  xlabel={leading coefficient on $u_{xx}$ (or $u_{tt}$)},
  ylabel={leading coefficient on $u_{yy}$ (or $u_{xx}$)},
  xmin=-1.4, xmax=1.4,
  ymin=-1.4, ymax=1.4,
  axis lines=middle,
  xtick=\empty, ytick=\empty,
  grid=none,
  font=\small,
  clip=false,
]
% Elliptic region: same-sign coefficients (quadrants I and III).
\addplot[fill=engblue!12, draw=none, forget plot] coordinates
  {(0,0) (1.4,0) (1.4,1.4) (0,1.4)} \closedcycle;
\addplot[fill=engblue!12, draw=none, forget plot] coordinates
  {(0,0) (-1.4,0) (-1.4,-1.4) (0,-1.4)} \closedcycle;
% Hyperbolic region: opposite-sign coefficients (quadrants II, IV).
\addplot[fill=engred!12, draw=none, forget plot] coordinates
  {(0,0) (1.4,0) (1.4,-1.4) (0,-1.4)} \closedcycle;
\addplot[fill=engred!12, draw=none, forget plot] coordinates
  {(0,0) (-1.4,0) (-1.4,1.4) (0,1.4)} \closedcycle;

% Marker: Laplace (elliptic) at (1,1).
\addplot[only marks, mark=*, engblue, mark size=2.5pt] coordinates {(1,1)};
\node[engblue, anchor=south west, font=\scriptsize] at (axis cs:1,1)
  {Laplace $u_{xx}+u_{yy}=0$};
% Marker: wave (hyperbolic) at (1,-1).
\addplot[only marks, mark=*, engred, mark size=2.5pt] coordinates {(1,-1)};
\node[engred, anchor=north west, font=\scriptsize] at (axis cs:1,-1)
  {wave $u_{tt}-c^{2}u_{xx}=0$};
% Marker: heat (parabolic) on the axis (one second derivative absent).
\addplot[only marks, mark=square*, engblack, mark size=2.5pt] coordinates {(1,0)};
\node[engblack, anchor=south, font=\scriptsize] at (axis cs:1,0.08)
  {heat $u_{t}-\alpha u_{xx}=0$ (parabolic)};

\node[engblue!60!engblack, font=\scriptsize, align=center]
  at (axis cs:-0.75,0.75) {elliptic\\ $B^{2}-AC<0$};
\node[engred!60!engblack, font=\scriptsize, align=center]
  at (axis cs:-0.75,-0.75) {hyperbolic\\ $B^{2}-AC>0$};
\end{axis}
\end{tikzpicture}
\caption[Classification of second-order linear PDEs]{Classification
of the second-order linear operator $A\,u_{xx} + 2B\,u_{xy} +
C\,u_{yy}$ by the sign of its discriminant $B^{2} - AC$, drawn for
the diagonal case $B = 0$ in the plane of the two leading
coefficients. Same-sign coefficients ($B^{2} - AC < 0$, blue)
give the elliptic type, of which Laplace's equation is the
prototype: steady state, no preferred direction, smoothing in
space. Opposite-sign coefficients ($B^{2} - AC > 0$, red) give the
hyperbolic type, of which the wave equation is the prototype:
finite-speed propagation along characteristics. The degenerate case
where one second derivative is absent ($AC = 0$, $B = 0$, so the
discriminant vanishes) gives the parabolic type, of which the heat
equation is the prototype: diffusive smoothing with an arrow of
time. The type fixes the qualitative behaviour, the appropriate
boundary and initial data, and the numerical schemes that
converge.}
\label{fig:vol02:ch12:classification}
\end{figure}
